\begin{table}[t] \centering \caption{Summary Statistics Indonesia}\label{tab:summary_stats_IDN_unb} \begin{threeparttable} \begin{tabular}{l cccc} \toprule
& All & Rural & Urban & Difference \\
 &  &  &  & $t$-test \\
\midrule
Location &  &   52.6\% &   47.4\% &  \\  \addlinespace \addlinespace
Log Consumption & 12.02 & 11.84 & 12.23 & -0.39*** \\
 & (0.79) & (0.77) & (0.76) &  \\  \addlinespace
Log Income & 14.88 & 14.64 & 15.15 & -0.50*** \\
 & (1.15) & (1.15) & (1.08) &  \\  \addlinespace
Female & 0.52 & 0.51 & 0.53 & -0.01*** \\
 & (0.50) & (0.50) & (0.50) &  \\  \addlinespace
Age (years) & 39.37 & 39.78 & 38.91 & 0.88*** \\
 & (14.95) & (15.36) & (14.48) &  \\  \addlinespace
Education (years) & 8.05 & 6.79 & 9.45 & -2.65*** \\
 & (4.65) & (4.45) & (4.46) &  \\  \addlinespace
Household Size & 4.89 & 4.78 & 5.00 & -0.22*** \\
 & (2.21) & (2.08) & (2.35) &  \\  \addlinespace


\cmidrule{2-5}
Observations & 93,038 & 48,972 & 44,066 &  \\  
Individuals & 29,716 &  &  &  \\
Non-switchers &   92.9\% &  &  &  \\  \addlinespace
\bottomrule \end{tabular} \begin{tablenotes}[flushleft] \footnotesize \item{Summary statistics for Indonesia for the unbalanced panel across all five waves. Source: IFLS. The table reports means and standard deviations (in parentheses) based on individual-year pairs. See section 3 for further details. All variables have the same number of observations, except for income, which is missing for some observations. Income has 61,300 observations.} \end{tablenotes} \end{threeparttable} \end{table}